\chapter{Classes of martingales and related processes}

The main reference for this part is \cite{he2019semimartingale}.

Notations for classes of processes:
\begin{itemize}
  \item $\mathcal{V}$: finite variation processes (Definition~\ref{def:HasFiniteVariation})
  \item $\mathcal{V}^+$: non-decreasing finite variation processes
  \item $\mathcal{A}$: adapted processes with integrable variation (Definition~\ref{def:HasIntegrableVariation})
  \item $\mathcal{A}^+$: non-decreasing adapted integrable processes
  \item $\mathcal{M}$: càdlàg martingales (TODO: \cite{he2019semimartingale} add uniform integrability)
  \item $\mathcal{M}^c$: continuous martingales
  \item $\mathcal{M}^2$: square-integrable martingales (Definition~\ref{def:IsSquareIntegrable})
  \item $\mathcal{C}_{\mathrm{loc}}$ for a class $\mathcal{C}$: processes that are locally in $\mathcal{C}$
\end{itemize}

Note: \cite{he2019semimartingale} use $\mathcal{W}$ for $\mathcal{M} \cap \mathcal{A}$.

\section{Jumps of a process}

\begin{definition}[Jumps of a process]\label{def:jump}
  \leanok
  \lean{Function.jump}
The jumps of a process $X : T \to \Omega \to E$ is the process $\Delta X : T \to \Omega \to E$ defined by $(t, \omega) \mapsto X_t(\omega) - X_{t^-}(\omega)$~.
\end{definition}


\begin{definition}[Jump part]\label{def:jumpPart}
  \uses{def:jump}
The jump part (or purely discontinuous part) of a process $X : T \to \Omega \to E$ is the process $T \to \Omega \to E$ defined by $(t, \omega) \mapsto \sum_{0 < s \le t} \Delta X_s(\omega)$~.
\end{definition}


\begin{definition}[Continuous part]\label{def:continuousPart}
  \uses{def:jumpPart}
The continuous part of a process $X : T \to \Omega \to E$ is the process $X^c : T \to \Omega \to E$ defined by $(t, \omega) \mapsto X_t(\omega) - \sum_{0 < s \le t} \Delta X_s(\omega)$~.
\end{definition}


\begin{definition}[Large jumps]\label{def:largeJump}
  \uses{def:jump}
The large jump part of a process $X : T \to \Omega \to E$ at level $\varepsilon$ is the process $X^{d, \varepsilon} : T \to \Omega \to E$ defined by $(t, \omega) \mapsto \sum_{0 < s \le t} \Delta X_s(\omega) \mathbb{1}_{\{\Vert \Delta X_s(\omega) \Vert > \varepsilon\}}$~.
\end{definition}


\section{Integrable variation}

\begin{definition}[Extended variation]\label{def:eVariationOn}
  \mathlibok
  \lean{eVariationOn}
The (extended-real-valued) variation of a function $f : T \to E$ on a set $s$ inside a linear order is the supremum of $\sum_i \mathrm{edist}(f(u_{i+1}, f(u_i)))$ over all finite increasing sequences $u : N \to T$ in $s$.
We denote it by $V_f(s)$~.
\end{definition}


\begin{definition}[Bounded variation]\label{def:BoundedVariationOn}
  \uses{def:eVariationOn}
  \mathlibok
  \lean{BoundedVariationOn}
A function $f : T \to E$ is of bounded variation on a set $s$ if its variation $V_f(s)$ is finite.
\end{definition}


\begin{definition}[Locally bounded variation]\label{def:LocallyBoundedVariationOn}
  \uses{def:BoundedVariationOn}
  \mathlibok
  \lean{LocallyBoundedVariationOn}
A function $f : T \to E$ is of locally bounded variation on a set $s$ if for every $a, b$ in $s$, the variation $V_f(s \cap [a, b])$ is finite.
\end{definition}


\begin{lemma}\label{lem:LocallyBoundedVariationOn.countable_not_continuous_at}
  \uses{def:LocallyBoundedVariationOn}
The set of points of discontinuity of a function of locally bounded variation is at most countable.
\end{lemma}

\begin{proof}
Mathlib has the result for monotone functions and has the fact that a function of locally bounded variation is the difference of two monotone functions.
\end{proof}


\begin{definition}[Finite variation process, $\mathcal{V}$]\label{def:HasFiniteVariation}
  \uses{def:LocallyBoundedVariationOn}
A process $X : T \to \Omega \to E$ is of finite variation if it is right-continuous and for every $\omega \in \Omega$, the path $t \mapsto X_t(\omega)$ is of locally bounded variation on $T$~.

We denote by $\mathcal{V}$ the class of finite variation processes.
\end{definition}


\begin{lemma}\label{lem:HasFiniteVariation.isCadlag}
  \uses{def:HasFiniteVariation, def:IsCadlag}
A finite variation process is càdlàg.
\end{lemma}

\begin{proof}
Right-continuity is by definition. Left limits exist because a function of locally bounded variation has left limits (that should be in Mathlib's file about locally bounded variation).
\end{proof}


\begin{definition}\label{def:variationProcess}
  \uses{def:eVariationOn, def:LocallyBoundedVariationOn}
The variation of a process $X : T \to \Omega \to E$ is the process $V_X : T \to \Omega \to \mathbb{R}$ defined by $(t, \omega) \mapsto V_{X(\omega)}([0,t])$~.
\end{definition}


\begin{lemma}\label{lem:monotone_variationProcess}
  \uses{def:variationProcess}
The variation process $V_X$ of a process $X$ is non-decreasing.
\end{lemma}

\begin{proof}

\end{proof}


\begin{lemma}\label{lem:rightContinuous_variationProcess}
  \uses{def:variationProcess}
The variation process $V_X$ of a right-continuous process $X$ is right-continuous.
\end{lemma}

\begin{proof}

\end{proof}


\begin{definition}[Integrable variation, $\mathcal{A}$]\label{def:HasIntegrableVariation}
  \uses{def:HasIntegrableSup, def:variationProcess}
We say that a stochastic process $X$ has integrable variation if its variation process $V_X$ (Definition~\ref{def:variationProcess}) is integrable (Definition~\ref{def:HasIntegrableSup}).

We denote by $\mathcal{A}$ the class of adapted processes with integrable variation.
\end{definition}

We denote by $\mathcal{A}^+$ the class of non-decreasing adapted integrable processes.

We denote by $\mathcal{V}^+$ the class of non-decreasing adapted processes.


\begin{lemma}\label{lem:HasIntegrableVariation.eq_sub_monotone}
  \uses{def:HasIntegrableVariation, def:HasIntegrableSup}
A process with integrable variation can be written as the difference of two non-decreasing adapted integrable processes.
\end{lemma}

\begin{proof}

\end{proof}


\begin{lemma}\label{lem:local_doobMeyer_integrable_variation}
  \uses{def:HasIntegrableVariation}
A process in $\mathcal{A}_{\mathrm{loc}}$ (locally integrable variation) can be written as a sum of a local martingale and a predictable locally integrable process.
\end{lemma}

\begin{proof}
  \uses{lem:HasIntegrableVariation.eq_sub_monotone, lem:local_doobMeyer_increasing}
We write the process as a difference of two non-decreasing locally integrable processes, and then apply Lemma~\ref{lem:local_doobMeyer_increasing} to each of them.
We take the difference of the two resulting local martingales and the difference of the two resulting predictable locally integrable processes.
\end{proof}


\begin{lemma}\label{lem:martingalePart_jumpPart}
  \uses{def:HasIntegrableVariation, def:Martingale}
Let $M \in \mathcal{M} \cap \mathcal{A}$ be a martingale with integrable variation. Then the jump part $A = \sum_{0 < s \le \cdot} \Delta M_s$ is an adapted process with integrable variation. Let $A_m$ be the martingale part of the Doob decomposition of $A$. Then $M = M_0 + A_m$. Finally, the predictable part of the decomposition of $A$ is continuous.
\end{lemma}

\begin{proof}

\end{proof}


\begin{lemma}\label{lem:Martingale.continuous_of_isStronglyPredictable_of_uniformIntegrable}
  \uses{def:Martingale, def:predictable}
A predictable uniformly integrable martingale is continuous.
\end{lemma}

\begin{proof}

\end{proof}


\section{Square integrable martingales}

In this section, $E$ denotes a complete normed space.

\begin{definition}[Square integrable martingales]\label{def:IsSquareIntegrable}
  \uses{def:Martingale}
  \leanok
  \lean{ProbabilityTheory.IsSquareIntegrable}
Let $T$ be a linear order with bottom element 0, on which we have a filtration $\mathcal{F}$ satisfying the usual conditions.
We say that a martingale $M : T \to \Omega \to E$ is square integrable if it is càdlàg and $\sup_{t \in T} \Vert M_t \Vert_{L^2} < \infty$ (Lean remark: use \texttt{eLpNorm (M t) 2}).
\end{definition}


\begin{lemma}\label{lem:IsSquareIntegrable.uniformIntegrable}
  \uses{def:IsSquareIntegrable}
  \leanok
  \lean{ProbabilityTheory.IsSquareIntegrable.uniformIntegrable}
  A square integrable martingale is uniformly integrable.
\end{lemma}

\begin{proof}
  \leanok
  \uses{lem:uniformIntegrable_of_bounded}
Apply Lemma~\ref{lem:uniformIntegrable_of_bounded}.
\end{proof}


\subsection{The Hilbert space of square integrable martingales}


\begin{lemma}\label{lem:IsSquareIntegrable.module}
  \uses{def:IsSquareIntegrable}
  \leanok
  \lean{ProbabilityTheory.IsSquareIntegrable.add, ProbabilityTheory.IsSquareIntegrable.smul}
If $M$ and $N$ are square integrable martingales and $a \in \mathbb{R}$, then $M + N$ and $a M$ are square integrable martingales.
\end{lemma}

\begin{proof}
\end{proof}


\begin{lemma}\label{lem:IsSquareIntegrable.submartingale_sq}
  \uses{def:IsSquareIntegrable, def:Submartingale}
  \leanok
  \lean{ProbabilityTheory.IsSquareIntegrable.submartingale_sq_norm}
If $M$ is a square integrable martingale, then $\Vert M \Vert^2$ is a submartingale.
\end{lemma}

\begin{proof}
  \uses{lem:Martingale.submartingale_convex_comp}
  \leanok
Apply Lemma~\ref{lem:Martingale.submartingale_convex_comp} with the convex function $f: x \mapsto \Vert x \Vert^2$~.
\end{proof}


\begin{lemma}\label{lem:IsSquareIntegrable.eLpNorm_two_mono}
  \uses{def:IsSquareIntegrable}
  \leanok
  \lean{ProbabilityTheory.IsSquareIntegrable.eLpNorm_mono}
For $M$ a square integrable martingale, the function $t \mapsto \Vert M_t \Vert_{L^2}$ is non-decreasing.
\end{lemma}

\begin{proof}
  \uses{lem:IsSquareIntegrable.submartingale_sq}
  \leanok
By Lemma~\ref{lem:IsSquareIntegrable.submartingale_sq}, $\Vert M_t \Vert^2$ is a submartingale.
Thus, for $s \le t$~,
\begin{align*}
  \Vert M_s \Vert_{L^2}^2
  &= \mathbb{E}[\Vert M_s \Vert^2]
  \\
  &\le \mathbb{E}[\Vert M_t \Vert^2]
  \\
  &= \Vert M_t \Vert_{L^2}^2
  \: .
\end{align*}
\end{proof}


\begin{lemma}\label{lem:IsSquareIntegrable.tendsto_limitProcess}
  \uses{def:IsSquareIntegrable, def:limitProcess}
  \leanok
  \lean{ProbabilityTheory.IsSquareIntegrable.ae_tendsto_limitProcess, ProbabilityTheory.IsSquareIntegrable.tendsto_eLpNorm_two_limitProcess}
For $M$ a square integrable martingale, we have $M_t \to M_\infty$ almost surely and in $L^2$ as $t \to \infty$~.
\end{lemma}

\begin{proof}
  \uses{thm:tendsto_limitProcess_of_cadlag}
TODO: use a martingale convergence theorem. Check whether Theorem~\ref{thm:tendsto_limitProcess_of_cadlag} is what we need.
\end{proof}


\begin{lemma}\label{lem:IsSquareIntegrable.limitProcess_add}
  \uses{def:IsSquareIntegrable, def:limitProcess}
  \leanok
  \lean{ProbabilityTheory.IsAESquareIntegrable.limitProcess_add}
For $M$ and $N$ two square integrable martingales, $(M + N)_\infty$ and $M_\infty + N_\infty$ are a.e. equal.
\end{lemma}

\begin{proof}
  \uses{lem:IsSquareIntegrable.tendsto_limitProcess}
Lemma~\ref{lem:limitProcess_ae_eq} with $g := M_\infty + N_\infty$.
\end{proof}


\begin{lemma}\label{lem:IsSquareIntegrable.limitProcess_sub}
  \uses{def:IsSquareIntegrable, def:limitProcess}
  \leanok
  \lean{ProbabilityTheory.IsAESquareIntegrable.limitProcess_sub}
For $M$ and $N$ two square integrable martingales, $(M - N)_\infty$ and $M_\infty - N_\infty$ are a.e. equal.
\end{lemma}

\begin{proof}
  \uses{lem:IsSquareIntegrable.tendsto_limitProcess}
Lemma~\ref{lem:limitProcess_ae_eq} with $g := M_\infty - N_\infty$.
\end{proof}


\begin{lemma}\label{lem:condExp_limitProcess}
  \uses{def:IsSquareIntegrable, def:limitProcess}
  \leanok
  \lean{ProbabilityTheory.IsSquareIntegrable.condExp_limitProcess_ae_eq}
  For $M$ a square integrable martingale and $t \in T$, we have that $P[X_\infty | \mathcal{F}_t]$ is a.e. equal to $X_t$.
\end{lemma}

\begin{proof}\leanok
  \uses{thm:condExp_limitProcess, lem:IsSquareIntegrable.uniformIntegrable}
  Apply Theorem~\ref{thm:condExp_limitProcess}.
\end{proof}


\begin{lemma}\label{lem:condExp_limitProcess_stopped}
  \uses{def:IsSquareIntegrable, def:limitProcess}
  \leanok
  \lean{ProbabilityTheory.IsSquareIntegrable.condExp_limitProcess_ae_eq'}
  For $M$ a square integrable martingale and $\tau$ a stopping time, we have that $P[X_\infty | \mathcal{F}_\tau]$ is a.e. equal to $X_\tau$.
\end{lemma}

\begin{proof}\leanok
  \uses{thm:condExp_limitProcess_stopped, lem:IsSquareIntegrable.uniformIntegrable}
  Apply Theorem~\ref{thm:condExp_limitProcess_stopped}.
\end{proof}


\begin{lemma}\label{lem:iSup_eLpNorm_le_eLpNorm_limitProcess}
  \uses{def:limitProcess}
  \leanok
  \lean{MeasureTheory.iSup_eLpNorm_le_eLpNorm_limitProcess}
For $M$ a càdlàg and uniformly integrable martingale,
$$\sup_{t \in T} \Vert M_t \Vert_2 \le \Vert M_\infty \Vert_2.$$
\end{lemma}

\begin{proof}
  \leanok
  \uses{lem:condExp_limitProcess}
By Lemma~\ref{lem:condExp_limitProcess}, We have that $P[M_\infty | \mathcal{F}_t] = M t$ almost surely, thus
$$\|M_t\|_2 = \|P[M_\infty | \mathcal{F}_t]\|_2 \le \|M_\infty\|_2,$$
where the last inequality is \href{https://leanprover-community.github.io/mathlib4_docs/Mathlib/MeasureTheory/Function/ConditionalExpectation/Real.html#MeasureTheory.eLpNorm_condExp_le_eLpNorm}{eLpNorm\_condExp\_le\_eLpNorm}.
\end{proof}


\begin{lemma}\label{lem:isSquareIntegrable_of_limitProcess}
  \uses{def:IsSquareIntegrable, def:limitProcess}
  \leanok
  \lean{ProbabilityTheory.isSquareIntegrable_of_limitProcess}
A càdlàg and uniformly integrable martingale $M$ is square integrable if and only if $M_\infty$ is in $\mathrm L^2$.
\end{lemma}

\begin{proof}
  \leanok
  \uses{lem:IsSquareIntegrable.iSup_eLpNorm_le_eLpNorm_limitProcess}
It is a càdlàg martingale by hypothesis, and the bound follows from Lemma~\ref{lem:IsSquareIntegrable.iSup_eLpNorm_le_eLpNorm_limitProcess}.
\end{proof}


\begin{lemma}\label{lem:IsSquareIntegrable.stoppedProcess}
  \uses{def:IsSquareIntegrable}
  \leanok
  \lean{ProbabilityTheory.IsSquareIntegrable.stoppedProcess}
If $M$ is a square integrable martingale and $\tau$ is a stopping time, then $M^\tau$ is a square integrable martingale.
\end{lemma}

\begin{proof}
  \leanok
  \uses{lem:isSquareIntegrable_of_limitProcess, lem:Martingale.classD_iff_uniformIntegrable, cor:doob_lp_norm_top, lem:IsSquareIntegrable.uniformIntegrable, lem:limitProcess_stoppedProcess}
We apply Lemma~\ref{lem:isSquareIntegrable_of_limitProcess}. First, $M^\tau$ is a martingale. Moreover, $M$ is uniformly integrable by Lemma~\ref{lem:IsSquareIntegrable.uniformIntegrable}, thus it is of class D by Lemma~\ref{lem:Martingale.classD_iff_uniformIntegrable}, thus $M^\tau$ is uniformly integrable. Finally, we have
\begin{align*}
  \|M^\tau_\infty\|_2^2 & = \|M_\tau\|^2 \\
  & \le \int (\sup_t \|M_t(\omega)\|)^2 P(\mathrm d\omega) \\
  & \le 4 \|M_\infty\|^2_2 \\
  & < \infty.
\end{align*}
The first equality follows from Lemma~\ref{lem:limitProcess_stoppedProcess}. The first inequality is obtained as follows: if $\tau(\omega) = \infty$, then $M_\tau = M_\infty$ and $M_t \longrightarrow M_\infty$ almost surely so $\|M_\infty(\omega)\| \le \sup_t \|M_t(\omega)\|$. Otherwise the inequality is obvious. The second inequality is Corollary~\ref{cor:doob_lp_norm_top}.
\end{proof}


\begin{lemma}\label{lem:IsSquareIntegrable.sup_eLpNorm_eq_eLpNorm_limitProcess}
  \uses{def:IsSquareIntegrable, def:limitProcess}
  \leanok
  \lean{ProbabilityTheory.IsSquareIntegrable.iSup_eLpNorm_eq_eLpNorm_limitProcess}
For $M$ a square integrable martingale,
\begin{align*}
  \sup_{t \in T} \Vert M_t \Vert_{L^2}
  &= \Vert M_\infty \Vert_{L^2}
  \: .
\end{align*}
\end{lemma}

\begin{proof}
  \uses{lem:IsSquareIntegrable.eLpNorm_two_mono, lem:IsSquareIntegrable.tendsto_limitProcess}

\end{proof}

\begin{lemma}\label{lem:tendsto_ae_condExp}
  \leanok
  \lean{ProbabilityTheory.tendsto_ae_condExp'}
  For $f : \Omega \to E$, we have almost surely that
  \begin{equation*}
    \lim_{t \to +\infty} P[f | \mathcal{F}_t] = P\left[f | \bigsqcup_t \mathcal{F}_t\right].
  \end{equation*}
\end{lemma}

\begin{proof}

\end{proof}


\begin{lemma}\label{lem:IsSquareIntegrable.integral_iSup_norm_rpow_rpow_inv_le_limitProcess}
  \uses{def:IsSquareIntegrable, def:limitProcess}
  \leanok
  \lean{ProbabilityTheory.IsSquareIntegrable.integral_iSup_norm_rpow_rpow_inv_le_limitProcess}
If $M$ is a square integrable random variable, then
$$\left(\int (\sup_t \|X_t(\omega)\|)^2\right)^{1/2} \mathrm d\omega \le 2 \|X_\infty\|_2.$$
\end{lemma}

\begin{proof}
  \leanok
  \uses{cor:doob_lp_norm_top, lem:IsSquareIntegrable.sup_eLpNorm_eq_eLpNorm_limitProcess}
Combine Corollary~\ref{cor:doob_lp_norm_top}, the fact that suprema commute with power, and Lemma~\ref{lem:IsSquareIntegrable.sup_eLpNorm_eq_eLpNorm_limitProcess}.
\end{proof}


\begin{lemma}\label{lem:memLp_two_stoppedValue}
  \uses{def:IsSquareIntegrable}
  \leanok
  \lean{ProbabilityTheory.IsAESquareIntegrable.memLp_two_stoppedValue'}
If $M$ is a square integrable martingale and $\tau$ is a stopping time, then $M_\tau$ is in $\mathrm L^2$.
\end{lemma}

\begin{proof}
  \leanok
  \uses{lem:limitProcess_stoppedProcess, lem:IsSquareIntegrable.stoppedProcess}
This follows from the fact that $M_\tau = M^\tau_\infty$ by Lemma~\ref{lem:limitProcess_stoppedProcess}, $M^\tau$ is square integrable by Lemma~\ref{lem:IsSquareIntegrable.stoppedProcess} and the stopped process of a square integrable martingale is in $\mathrm L^2$.
\end{proof}


\begin{definition}\label{def:IsPurelyDiscontinuous}
  \uses{def:IsSquareIntegrable}
  \leanok
  \lean{ProbabilityTheory.IsPurelyDiscontinuous}
A stochastic process $X$ is purely discontinuous if it is a square integrable martingale such that for any continuous square integrable martingale $Y$,
$$P[\langle X_\infty, Y_\infty\rangle] = 0.$$
\end{definition}


\begin{definition}\label{def:L2Martingales}
  \uses{def:IsSquareIntegrable}
  \leanok
  \lean{ProbabilityTheory.SquareIntegrable}
We denote by $\mathcal{M}^2(E)$ or simply $\mathcal{M}^2$ the space of equivalence classes with respect to indistinguishability of square integrable martingales $T \to \Omega \to E$~.
\end{definition}


\begin{lemma}\label{lem:L2Martingales.module}
  \uses{def:L2Martingales}
  \leanok
  \lean{ProbabilityTheory.SquareIntegrable}
The space $\mathcal{M}^2(E)$ is a real vector space.
\end{lemma}

\begin{proof}
  \uses{lem:IsSquareIntegrable.module, lem:limitProcess_congr, lem:IsSquareIntegrable.limitProcess_add, lem:limitProcess_smul}
\end{proof}


\begin{definition}\label{def:L2Martingales.norm}
  \uses{def:L2Martingales, def:limitProcess}
  \leanok
  \lean{ProbabilityTheory.instNormedAddCommGroupSquareIntegrable}
We define a norm on $\mathcal{M}^2$ by
\begin{align*}
  \Vert M \Vert = \Vert M_\infty \Vert_{L^2}
  \: .
\end{align*}
\end{definition}


\begin{lemma}\label{lem:L2Martingales.norm_eq_zero}
  \uses{def:L2Martingales.norm}
  \leanok
  \lean{ProbabilityTheory.SquareIntegrable.injective_toL2}
For $M \in \mathcal{M}^2(E)$, $\Vert M \Vert = 0$ if and only if $M = 0$.
\end{lemma}

\begin{proof}
  \uses{lem:IsSquareIntegrable.sup_eLpNorm_eq_eLpNorm_limitProcess}
By Lemma~\ref{lem:IsSquareIntegrable.sup_eLpNorm_eq_eLpNorm_limitProcess}, $\Vert M \Vert = 0$ if and only if for all $t \in T$, $\Vert M_t \Vert_{L^2} = 0$~. Because $M$ is càdlàg, this implies that $M$ is indistinguishable from $0$.
\end{proof}


\begin{definition}\label{def:L2Martingales.inner}
  \uses{def:L2Martingales.norm}
  \leanok
  \lean{ProbabilityTheory.instInnerProductSpaceRealSquareIntegrable}
We define an inner product on $\mathcal{M}^2$ by
\begin{align*}
  \langle M, N \rangle_{\mathcal{M}^2} = \mathbb{E}[M_\infty N_\infty]
  \: .
\end{align*}
\end{definition}


\begin{theorem}\label{thm:hilbertSpace_L2Martingales}
  \uses{def:L2Martingales.inner, lem:L2Martingales.norm_eq_zero}
  \leanok
  \lean{ProbabilityTheory.instCompleteSpaceSquareIntegrableOfOrderTopology}
If $T$ is separable then the space $\mathcal{M}^2$ is a Hilbert space.
\end{theorem}

\begin{proof}
  \uses{lem:L2Martingales.module, lem:exists_cadlag_mod_of_local_mg}
  We show that the map
  \begin{align*}
    \phi : \mathcal{M}^2 & \to \left\{f \in L^2(E, P) | f \text{ is } \left(\bigsqcup_t, \mathcal{F}_t\right)\text{-strongly measurable}\right\} \\
    M & \mapsto M_\infty
  \end{align*}
  is a bijective isometry. It is an isometry by Definition~\ref{def:L2Martingales.norm}. We define the inverse $\psi$ as the map which to $f$ associates a càdlàg martingale that is a modification of $t \mapsto P[f | \mathcal{F}_t]$, given by Lemma~\ref{lem:exists_cadlag_mod_of_local_mg}.

  Let $M$ be a square integrable martingale. Then for any $t$, $M_t$ is a.e. equal to $P[M_\infty | \mathcal{F}_t]$ by Lemma~\ref{lem:condExp_limitProcess}, thus $M$ is a càdlàg modification of $t \mapsto P[f | \mathcal{F}_t]$, thus it is indistinguishable from $\psi(M_\infty)$.

  Let $f \in L^2(E, P)$ be a.e. strongly measurable with respect to $\bigsqcup_t \mathcal{F}_t$. Then by Lemma~\ref{lem:tendsto_ae_condExp} we know that almost surely $P[f | \mathcal{F}_t]$ converges to $f$. Because $T$ is separable we can consider $(t_n)$ a sequence of points that tends to infinity. Because $\psi(f)$ and $P[f | \mathcal{F}_t]$ are modifications of each other, we have that almost surely, for all $n \in \mathcal{N}$, $\psi(f)_{t_n} = P[f | \mathcal{F}_{t_n}]$. But we also know that $\psi(f)_{t_n}$ converges to $\psi(f)_\infty$ so we conclude that $\psi(f)_\infty$ is a.e. equal to f.
\end{proof}





\subsection{Elementary stochastic integrals}

\begin{lemma}\label{lem:eLpNorm_elemStochIntegralBilin_le}
  \uses{def:IsSquareIntegrable, def:elemStochIntegralBilin}
For $V \in \mathcal{E}_{T, F}$ bounded by a constant $D$, $M \in \mathcal{M}^2(E)$ and a continuous bilinear map $B: E \times F \to G$,
\begin{align*}
  \Vert (V \bullet_B M)_t \Vert_{L^2}
  \le 2 D \: \Vert B \Vert \: \sup_t \Vert M_t \Vert_{L^2}
\end{align*}
TODO: this can be improved to $D \: \Vert B \Vert \: \Vert M_t \Vert_{L^2}$?
\end{lemma}

\begin{proof}
Let $C$ be a bound on $\Vert M_t \Vert_{L^2}$ for all $t \in T$~.
Let $(s_k < t_k)_{k \in \{1, ..., n\}}$ and $\eta_k$ be the intervals and random variables defining $V$~.
Let $D$ be a bound on $\Vert\eta_k\Vert$.
Then, for all $t$~,
\begin{align*}
  \Vert (V \bullet_B M)_t \Vert_{L^2}
  &\le \sum_{k=1}^n \Vert B(M^t_{t_k} - M^t_{s_k}, \eta_k) \Vert_{L^2}
  \: .
\end{align*}
Since only at most one term of that sum is non-zero for each fixed $t$~, we can bound the sum by the maximum of its terms.
It suffices then to bound each term of that sum.

TODO: here we supposed that the intervals of the simple process are disjoint. Check with our Lean def.

For each $k$~,
\begin{align*}
  \Vert B(M^t_{t_k} - M^t_{s_k}, \eta_k) \Vert_{L^2}
  &\le \left\Vert \Vert B \Vert \: \Vert M^t_{t_k} - M^t_{s_k} \Vert  \: \Vert \eta_k \Vert \right\Vert_{L^2}
  \\
  &\le \Vert B \Vert \: \Vert M^t_{t_k} - M^t_{s_k} \Vert_{L^2} \: D
  \\
  &\le 2 \Vert B \Vert \: C \: D
  \: .
\end{align*}
\end{proof}


\begin{lemma}\label{lem:isSquareIntegrable_elemStochIntegralBilin}
  \uses{def:IsSquareIntegrable, def:elemStochIntegralBilin}
For $V \in \mathcal{E}_{T, F}$, $M \in \mathcal{M}^2(E)$ and a continuous bilinear map $B: E \times F \to G$, the elementary stochastic integral $V \bullet_B M$ is in $\mathcal{M}^2(G)$.
\end{lemma}

\begin{proof}
  \uses{lem:cadlag_elemStochIntegralBilin, lem:Martingale.elemStochIntegral, lem:eLpNorm_elemStochIntegralBilin_le}
By Lemma~\ref{lem:cadlag_elemStochIntegralBilin}, $V \bullet_B M$ is càdlàg, and we know that it is a martingale by Lemma~\ref{lem:Martingale.elemStochIntegral}~.
It remains to show that $\sup_{t \in T} \Vert (V \bullet_B M)_t \Vert_{L^2} < \infty$~.
By Lemma~\ref{lem:eLpNorm_elemStochIntegralBilin_le}, this supremum is bounded by $2 D \Vert B \Vert \sup_{t \in T} \Vert M_t \Vert_{L^2}$, which is finite since $M \in \mathcal{M}^2(E)$ and $V$ is bounded.
\end{proof}


\begin{definition}\label{def:continuousSquareIntegrable}
  \uses{def:L2Martingales}
  \leanok
  \lean{ProbabilityTheory.continuousSquareIntegrable}
We denote by $\mathcal{M}^{2,c}$ the submodule of $\mathcal{M}^2$ made of continuous square integrable martingales.
\end{definition}


\begin{lemma}\label{lem:exists_subsequence_ae_tendsto_uniformly}
  \uses{def:IsSquareIntegrable}
  \leanok
  \lean{ProbabilityTheory.exists_subsequence_ae_tendsto_uniformly}
Let $(M^{(n)})$ be a sequence of square integrable martingales and $N$ be a square integrable martingale such that $\|M^{(n)}_\infty - N_\infty\|_2 \longrightarrow_{n \to +\infty} 0$. Then there exists $\phi : \mathbb{N} \to \mathbb{N}$ increasing such that, almost surely, $M^{(\phi(n))}_\cdot(\omega)$ converges uniformly to $N_\cdot(\omega)$ as $n$ goes to infinity.
\end{lemma}

\begin{proof}
  \leanok
  \uses{lem:IsSquareIntegrable.integral_iSup_norm_rpow_rpow_inv_le_limitProcess}
We can find $\phi$ increasing such that $\|M^{(\phi(n))}_\infty - N_\infty\|_2 \le 1/2^n$. Then,
\begin{align*}
  P[\sum_n \sup_t \|M^{(\phi(n))}_t -  N_t\|] & = \sum_n P[\sup_t \|M^{(\phi(n))}_t -  N_t\|] \\
  & \le \sqrt{P(\Omega)} \sum_n P[(\sup_t \|M^{(\phi(n))}_t -  N_t\|)^2]^{1/2} \\
  & \le 2 \sqrt{P(\Omega)} \|M^{(\phi(n))}_\infty - N_\infty\|_2 < \infty,
\end{align*}
Where the second inequality is Lemma~\ref{lem:IsSquareIntegrable.integral_iSup_norm_rpow_rpow_inv_le_limitProcess}. In particular, almost surely, $\sum_n \sup_t \|M^{(\phi(n))}_t -  N_t\| < \infty$ and thus $\sup_t \|M^{(\phi(n))}_t -  N_t\| \longrightarrow_{n \to +\infty} 0$, which concludes the proof.
\end{proof}


\begin{lemma}\label{lem:isClosed_continuousSquareIntegrable}
  \uses{def:continuousSquareIntegrable}
  \leanok
  \lean{ProbabilityTheory.instIsClosedSquareIntegrableCoeSubmoduleRealContinuousSquareIntegrableOfIsCompleteOfOrderTopology}
The submodule $\mathcal{M}^{2,c}$ is closed in the Hilbert space $\mathcal{M}^2$.
\end{lemma}

\begin{proof}
  \leanok
  \uses{lem:exists_subsequence_ae_tendsto_uniformly}
Take $M^{(n)}$ a sequence of continuous square integrable martingales that converges to $N$ in $\mathcal{M}^2$. This means that $\|M^{(n)}_\infty - N_\infty\|$ goes to zero as $n$ goes to infinity, which means that there exists $\phi$ increasing such that, almost surely, $M^{(\phi(n))}_\cdot(\omega)$ converges uniformly to $N_\cdot(\omega)$ as $n$ goes to infinity (Lemma~\ref{lem:exists_subsequence_ae_tendsto_uniformly}). As each $M^{(\phi(n))}(\omega)$ is continuous, $N(\omega)$ is also continuous.
\end{proof}


\begin{definition}\label{def:continuousPartSquareIntegrable}
  \uses{def:continuousSquareIntegrable, lem:isClosed_continuousSquareIntegrable}
  \leanok
  \lean{ProbabilityTheory.continuousPart}
If $M$ is a square integrable martingale, we define its \emph{continuous part} as the orthogonal projection of $M$ onto the closed subspace $\mathcal{M}^{2,c}$. We denote it by $M^c$.
\end{definition}


\begin{lemma}\label{lem:continuous_continuousPart}
  \uses{def:continuousPartSquareIntegrable}
  \leanok
  \lean{ProbabilityTheory.continuous_continuousPart}
The continuous part of a square integrable martingale is continuous.
\end{lemma}

\begin{proof}
  \leanok
  \uses{def:continuousPartSquareIntegrable}
By definition of the orthogonal projection.
\end{proof}


\begin{lemma}\label{lem:isSquareIntegrable_continuousPart}
  \uses{def:continuousSquareIntegrable}
  \leanok
  \lean{ProbabilityTheory.isSquareIntegrable_continuousPart}
The continuous part of a square integrable martingale is a square integrable martingale.
\end{lemma}

\begin{proof}
  \leanok
  \uses{def:continuousPartSquareIntegrable}
By definition of the orthogonal projection.
\end{proof}


\begin{definition}\label{def:discontinuousPart}
  \uses{def:continuousPartSquareIntegrable}
  \leanok
  \lean{ProbabilityTheory.discontinuousPart}
If $M$ is a square integrable martingale, we define its \emph{discontinuous part} by $M^d := M - M^c$.
\end{definition}


\begin{lemma}\label{lem:isSquareIntegrable_discontinuousPart}
  \uses{def:continuousSquareIntegrable}
  \leanok
  \lean{ProbabilityTheory.isSquareIntegrable_continuousPart}
The discontinuous part of a square integrable martingale is a square integrable martingale.
\end{lemma}

\begin{proof}
  \leanok
  \uses{def:discontinuousPart, lem:isSquareIntegrable_continuousPart}
By definition.
\end{proof}


\begin{definition}\label{def:discontinuousSquareIntegrable}
  \uses{def:continuousSquareIntegrable}
  \leanok
  \lean{ProbabilityTheory.discontinuousSquareIntegrable}
We denote by $\mathcal{M}^{2,d}$ the orthogonal of $\mathcal{M}^{2,c}$ and call its elements purely discontinuous square integrable martingales.
\end{definition}


\begin{lemma}\label{lem:mem_discontinuousSquareIntegrable}
  \uses{def:IsPurelyDiscontinuous, def:discontinuousSquareIntegrable}
  \leanok
  \lean{ProbabilityTheory.mem_discontinuousSquareIntegrable}
If $M$ is a purely discontinuous martingale in the sense of Definition~\ref{def:IsPurelyDiscontinuous}, then it belongs to $\mathcal{M}^{2,d}$.
\end{lemma}

\begin{proof}
  \leanok
By definition of the orthogonal projection.
\end{proof}


\begin{lemma}\label{lem:isPurelyDiscontinuous_discontinuousPart}
  \uses{def:IsPurelyDiscontinuous, def:discontinuousPart}
  \leanok
  \lean{ProbabilityTheory.isPurelyDiscontinuous_discontinuousPart}
The discontinuous part of square integrable martingale is purely discontinuous.
\end{lemma}

\begin{proof}
  \leanok
By definition of the orthogonal projection.
\end{proof}


\begin{lemma}\label{lem:unique_continuousPart}
  \uses{def:continuousPartSquareIntegrable, def:discontinuousPart}
  \leanok
  \lean{ProbabilityTheory.indist_continuousPart, ProbabilityTheory.indist_discontinuousPart}
If $M$ is a square integrable martingale, then the decomposition $M = M^c + M^d$ as the sum of a continuous and a purely discontinuous square integrable martingales is unique up to indistinguishability.
\end{lemma}

\begin{proof}
  \leanok
A closed subspace is the complement of its orthogonal subspace.
\end{proof}


\begin{lemma}\label{lem:IsPurelyDiscontinuous.stoppedProcess}
  \uses{def:IsPurelyDiscontinuous}
  \leanok
  \lean{ProbabilityTheory.IsPurelyDiscontinuous.stoppedProcess}
If $M$ is a purely discontinuous martingale and $\tau$ is a stopping time, then $M^\tau$ is a purely discontinuous martingale.
\end{lemma}

\begin{proof}
  \leanok
  \uses{lem:IsSquareIntegrable.stoppedProcess, lem:limitProcess_stoppedProcess, lem:condExp_limitProcess_stopped}
By Lemma~\ref{lem:IsSquareIntegrable.stoppedProcess} $M^\tau$ is a square integrable martingale. Let $N$ be a continuous square integrable martingale. Then again by Lemma~\ref{lem:IsSquareIntegrable.stoppedProcess}, $N^\tau$ is a continuous square integrable martingale. We thus have
\begin{align*}
  P[\langle M^\tau_\infty, N_\infty\rangle] & = P[\langle M_\tau, N_\infty\rangle] \\
  & = P[\langle M_\tau, P[M_\infty | \mathcal{F}_\tau]\rangle] \\
  & = P[\langle M_\tau, N_\tau\rangle] \\
  & = P[\langle M_\infty, N_\tau\rangle] \\
  & = 0.
\end{align*}
The first equality is Lemma~\ref{lem:limitProcess_stoppedProcess}, the second is the pull-out property of the conditional expectation, the third one is Lemma~\ref{lem:condExp_limitProcess_stopped}, the fourth one is the same as before and the last one is the fact that $M$ is purely discontinuous while $N^\tau$ is continuous.
\end{proof}


\begin{lemma}\label{lem:continuousPart_stoppedProcess}
  \uses{def:continuousPartSquareIntegrable}
  \leanok
  \lean{ProbabilityTheory.continuousPart_stoppedProcess, ProbabilityTheory.discontinuousPart_stoppedProcess}
If $M$ is a square integrable martingale and $\tau$ is a stopping time, then $(M^\tau)^c = (M^c)^\tau$ and $(M^\tau)^d = (M^d)^\tau$.
\end{lemma}

\begin{proof}
  \leanok
  \uses{lem:unique_continuousPart, lem:IsPurelyDiscontinuous.stoppedProcess}
We can write $M^\tau = (M^c + M^d)^\tau = (M^c)^\tau + (M^d)^\tau$. We conclude by Lemma~\ref{lem:unique_continuousPart} thanks to Lemma~\ref{lem:IsPurelyDiscontinuous.stoppedProcess}.
\end{proof}



\begin{lemma}\label{lem:inner_elemStochIntegral}
  \uses{def:IsSquareIntegrable, def:elemStochIntegralBilin, lem:isSquareIntegrable_elemStochIntegralBilin}
For $V \in \mathcal{E}_{T, \mathbb{R}}$ and $M, N \in \mathcal{M}^2$, we have
\begin{align*}
  \langle V \bullet_{\mathbb{R}} M, N \rangle_{\mathcal{M}^2}
  &= V \bullet_{\mathbb{R}} \langle M, N \rangle_{\mathcal{M}^2}
  \: .
\end{align*}
\end{lemma}

\begin{proof}

\end{proof}



\section{Locally square integrable martingales}

\subsection{Definition and basic properties}

\begin{definition}[Locally square-integrable martingales]\label{def:IsLocallySquareIntegrable}
  \uses{def:IsSquareIntegrable, def:locally}
  \leanok
  \lean{ProbabilityTheory.IsLocallySquareIntegrable}
A process is locally square-integrable if it locally satisfies the square-integrable martingale property.
We denote that class of processes by $\mathcal{M}^2_{\mathrm{loc}}$~.
\end{definition}


\begin{lemma}\label{lem:IsSquareIntegrable.isLocallySquareIntegrable}
  \uses{def:IsSquareIntegrable, def:IsLocallySquareIntegrable}
  \leanok
  \lean{ProbabilityTheory.IsSquareIntegrable.isLocallySquareIntegrable}
Every square-integrable martingale is locally square-integrable: $\mathcal{M}^2 \subseteq \mathcal{M}^2_{\mathrm{loc}}$~.
\end{lemma}

\begin{proof}\leanok
  \uses{lem:implies_locally}
This follows from Lemma~\ref{lem:implies_locally}.
\end{proof}


\begin{lemma}\label{lem:IsLocallySquareIntegrable.isLocalSubmartingale_sq_norm}
  \uses{def:IsLocallySquareIntegrable, def:IsLocalSubmartingale}
  \leanok
  \lean{ProbabilityTheory.IsLocallySquareIntegrable.isLocalSubmartingale_sq_norm}
If $M \in \mathcal{M}^2_{\mathrm{loc}}$, then $\Vert M \Vert^2$ is a càdlàg local submartingale.
\end{lemma}

\begin{proof}\leanok
  \uses{lem:IsSquareIntegrable.submartingale_sq}

\end{proof}


\begin{lemma}\label{lem:IsLocalMartingale.isLocallySquareIntegrable_of_continuous}
  \uses{def:IsLocalMartingale, def:IsLocallySquareIntegrable}
A continuous local martingale is locally square-integrable: $\mathcal{M}^c_{\mathrm{loc}} \subseteq \mathcal{M}^2_{\mathrm{loc}}$~.
\end{lemma}

\begin{proof}

\end{proof}



\subsection{Predictable quadratic variation}


\begin{definition}[Predictable quadratic variation]\label{def:quadraticVariation}
  \uses{def:IsLocallySquareIntegrable, def:IsCadlag, thm:local_doobMeyer,
    lem:IsLocallySquareIntegrable.isLocalSubmartingale_sq_norm}
  \leanok
  \lean{ProbabilityTheory.predQuadVariation}
For $M \in \mathcal{M}^2_{\mathrm{loc}}$ with càdlàg paths, the predictable
quadratic variation of $M$ is defined as the predictable part of the Doob-Meyer decomposition of the local submartingale $\Vert M \Vert^2$~.
We denote it by $\langle M \rangle$~.
\end{definition}


\begin{lemma}\label{lem:predictable_quadraticVariation}
  \uses{def:quadraticVariation}
  \leanok
  \lean{ProbabilityTheory.isStronglyPredictable_predQuadVariation}
The predictable quadratic variation $\langle M \rangle$ of $M \in \mathcal{M}^2_{\mathrm{loc}}$ is a predictable process.
\end{lemma}

\begin{proof}\leanok
  \uses{thm:local_doobMeyer}

\end{proof}


\begin{lemma}\label{lem:cadlag_quadraticVariation}
  \uses{def:quadraticVariation}
  \leanok
  \lean{ProbabilityTheory.isCadlag_predQuadVariation}
The predictable quadratic variation $\langle M \rangle$ of $M \in \mathcal{M}^2_{\mathrm{loc}}$ is càdlàg.
\end{lemma}

\begin{proof}\leanok
  \uses{thm:local_doobMeyer}

\end{proof}


\begin{lemma}\label{lem:locallyIntegrable_quadraticVariation}
  \uses{def:quadraticVariation}
  \leanok
  \lean{ProbabilityTheory.hasLocallyIntegrableSup_predQuadVariation}
The predictable quadratic variation $\langle M \rangle$ of $M \in \mathcal{M}^2_{\mathrm{loc}}$ is locally integrable.
\end{lemma}

\begin{proof}\leanok
  \uses{thm:local_doobMeyer}

\end{proof}


\begin{lemma}\label{lem:quadraticVariation_zero}
  \uses{def:quadraticVariation}
$\langle M \rangle_0 = 0$~.
\end{lemma}

\begin{proof}
  \uses{thm:local_doobMeyer}

\end{proof}


\begin{lemma}\label{lem:monotone_quadraticVariation}
  \uses{def:quadraticVariation}
  \leanok
  \lean{ProbabilityTheory.monotone_predQuadVariation}
The predictable quadratic variation $\langle M \rangle$ of $M \in \mathcal{M}^2_{\mathrm{loc}}$ is non-decreasing.
\end{lemma}

\begin{proof}\leanok
  \uses{thm:local_doobMeyer}

\end{proof}


\begin{lemma}\label{lem:local_martingale_sub_quadraticVariation}
  \uses{def:quadraticVariation}
For $M \in \mathcal{M}^2_{\mathrm{loc}}$, the process $\Vert M_t \Vert^2 - \langle M \rangle_t$ is a local martingale.
\end{lemma}

\begin{proof}
  \uses{thm:local_doobMeyer}

\end{proof}


\begin{definition}[Predictable covariation]\label{def:covariation}
  \uses{def:IsLocallySquareIntegrable, def:quadraticVariation}
  \leanok
  \lean{ProbabilityTheory.predQuadCovariation}
For $M, N \in \mathcal{M}^2_{\mathrm{loc}}$, the predictable covariation $\langle M, N \rangle$ is a stochastic process defined by polarization of the predictable quadratic variation:
\begin{align*}
  \langle M, N \rangle_t = \frac{1}{2}\left(\langle M+N \rangle_t - \langle M \rangle_t - \langle N \rangle_t \right)
  \: .
\end{align*}
\end{definition}


\begin{lemma}\label{lem:predictable_covariation}
  \uses{def:covariation}
  \leanok
  \lean{ProbabilityTheory.isStronglyPredictable_predQuadCovariation}
The predictable covariation $\langle M, N \rangle$ of $M, N \in \mathcal{M}^2_{\mathrm{loc}}$ is a predictable process.
\end{lemma}

\begin{proof}\leanok
  \uses{lem:predictable_quadraticVariation}

\end{proof}


\begin{lemma}\label{lem:cadlag_covariation}
  \uses{def:covariation}
  \leanok
  \lean{ProbabilityTheory.isCadlag_predQuadCovariation}
The predictable covariation $\langle M, N \rangle$ of $M, N \in \mathcal{M}^2_{\mathrm{loc}}$ is càdlàg.
\end{lemma}

\begin{proof}\leanok
  \uses{lem:cadlag_quadraticVariation}

\end{proof}


\begin{lemma}\label{lem:covariation_zero}
  \uses{def:covariation}
$\langle M, N \rangle_0 = 0$~.
\end{lemma}

\begin{proof}
  \uses{lem:quadraticVariation_zero}

\end{proof}


\begin{lemma}\label{lem:local_martingale_sub_covariation}
  \uses{def:covariation}
For $M, N \in \mathcal{M}^2_{\mathrm{loc}}$, the process $\langle M_t, N_t \rangle_E - \langle M, N \rangle_t$ is a local martingale.
\end{lemma}

\begin{proof}
  \uses{lem:local_martingale_sub_quadraticVariation}
\begin{align*}
  &\langle M_t, N_t \rangle_E - \langle M, N \rangle_t
  \\
  &= \frac{1}{4}\left( \left(\Vert M_t + N_t \Vert^2 - \langle M+N \rangle_t\right) - \left(\Vert M_t - N_t \Vert^2 - \langle M-N \rangle_t\right) \right)
  \: .
\end{align*}
The two differences are local martingales by Lemma~\ref{lem:local_martingale_sub_quadraticVariation}, so their linear combination is also a local martingale.
\end{proof}


\begin{lemma}\label{lem:covariation_eq_inner}
  \uses{def:covariation, def:IsSquareIntegrable}
Let $M$ and $N$ be square integrable martingales. Then
\begin{align*}
  \mathbb{E}\left[\langle M,N \rangle_\infty\right] = \langle M - M_0, N - N_0 \rangle_{\mathcal{M}^2}
  \: .
\end{align*}
\end{lemma}

\begin{proof}

\end{proof}


% TODO: move elsewhere?
\begin{lemma}\label{lem:quadraticVariation_brownian}
  \uses{def:brownian, def:quadraticVariation}
Let $B$ be a standard Brownian motion. Then the quadratic variation of $B$ is given by $\langle B \rangle_t = t$~.
\end{lemma}

\begin{proof}

\end{proof}




\section{Local martingales}

\begin{lemma}\label{lem:IsLocalMartingale.locally_hasIntegrableVariation_largeJumps}
  \uses{def:IsLocalMartingale, def:HasIntegrableVariation, def:largeJump}
The large jump process of a local martingale is a process with locally integrable variation (it's in $(\mathcal{A} \cap \mathcal{M})_{\mathrm{loc}}$).
\end{lemma}

\begin{proof}

\end{proof}


\begin{theorem}\label{thm:local_martingale_decomposition}
  \uses{def:IsLocalMartingale}
Let $M$ be a local martingale.
Then for any $\varepsilon > 0$, $M$ can be decomposed as $M = M_0 + U + V$, where $U$ is a locally bounded martingale (local version of both bounded and martingale) with $\vert \Delta U \vert \le \varepsilon$ and $U_0 = 0$ and $V$ is a local martingale with locally integrable variation ($V \in (\mathcal{M} \cap \mathcal{A})_{\mathrm{loc}}$) and $V_0 = 0$.
\end{theorem}

\begin{proof}
  \uses{lem:IsLocalMartingale.locally_hasIntegrableVariation_largeJumps}
See  \cite{he2019semimartingale}, 7.17
\end{proof}

Remark:
$U$ is locally bounded, hence locally square integrable, so we can use the integration machinery for those to define an integral.
$V$ has locally integrable variation so we can integrate it with Stieltjes integrals.


\section{Semimartingales}


\begin{definition}[Semimartingale]\label{def:IsSemimartingale}
  \uses{def:IsLocalMartingale, def:HasFiniteVariation, def:adapted}
A process $X$ is a semimartingale if it can be decomposed as $X = M + A$, where $M$ is a local martingale ($M \in \mathcal{M}_{\mathrm{loc}}$) and $A$ is an adapted process with finite variation.
\end{definition}

TODO: it's adapted and cadlag.


\begin{lemma}\label{lem:IsSemimartingale.decomposition}
  \uses{def:IsSemimartingale}
A semimartingale $X$ can be decomposed as $X = M + A$, where $M$ is a locally bounded martingale with bounded jumps and $A$ is an adapted process with finite variation.
\end{lemma}

\begin{proof}
  \uses{thm:local_martingale_decomposition}
Decompose $X = M + A$ as in Definition~\ref{def:IsSemimartingale}.
Then decompose $M$ as in Theorem~\ref{thm:local_martingale_decomposition} with $\varepsilon = 1$.
Then $M = M_0 + U + V$ with $U$ a locally bounded martingale with $\vert \Delta U \vert \le 1$ and $V$ a local martingale with locally integrable variation.
Then $X = (M_0 + U) + (A + V)$, with $U$ a locally bounded martingale with bounded jumps and $A + V$ an adapted process of finite variation (since $V$ has locally integrable variation, hence finite variation).
\end{proof}


\begin{definition}[Continuous martingale part]\label{def:IsSemimartingale.continuousMartingalePart}
  \uses{def:IsSemimartingale, lem:IsSemimartingale.decomposition}
TODO. Denoted by $X^c$.
\end{definition}


\begin{definition}[Quadratic covariation of semimartingales]\label{def:quadCovariation}
  \uses{def:IsSemimartingale, def:quadraticVariation, def:IsSemimartingale.continuousMartingalePart}
Let $X$ and $Y$ be semimartingales. Their quadratic covariation $[X, Y]$ is defined as
\begin{align*}
  [X, Y]_t = X_0 Y_0 + \langle X^c, Y^c \rangle_t + \sum_{0 < s \le t} \Delta X_s \Delta Y_s
  \: .
\end{align*}
\end{definition}

TODO: this is an adapted process with finite variation (it's in $\mathcal{V}$).


\begin{definition}[Quadratic variation of a semimartingale]\label{def:quadVariation}
  \uses{def:quadCovariation}
The quadratic variation of a semimartingale $X$ is defined as $[X] = [X, X]$~.
\end{definition}

TODO: $[X]$ is an adapted increasing process.


\begin{definition}[Special semimartingale]\label{def:IsSpecialSemimartingale}
  \uses{def:IsSemimartingale, def:HasIntegrableVariation}
A semimartingale $X$ is a special semimartingale if it can be decomposed as $X = M + A$, where $M$ is a local martingale and $A$ is an adapted process with locally integrable variation.
\end{definition}


\begin{theorem}\label{lem:IsSpecialSemimartingale.decomposition}
  \uses{def:IsSpecialSemimartingale}
A special semimartingale $X$ can be decomposed as $X = M + A$, where $M$ is a local martingale and $A$ is a predictable process with finite variation and $A_0 = 0$.
\end{theorem}

\begin{proof}

\end{proof}
